% Options for packages loaded elsewhere
\PassOptionsToPackage{unicode}{hyperref}
\PassOptionsToPackage{hyphens}{url}
\documentclass[
]{article}
\usepackage{xcolor}
\usepackage[margin=1in]{geometry}
\usepackage{amsmath,amssymb}
\setcounter{secnumdepth}{-\maxdimen} % remove section numbering
\usepackage{iftex}
\ifPDFTeX
  \usepackage[T1]{fontenc}
  \usepackage[utf8]{inputenc}
  \usepackage{textcomp} % provide euro and other symbols
\else % if luatex or xetex
  \usepackage{unicode-math} % this also loads fontspec
  \defaultfontfeatures{Scale=MatchLowercase}
  \defaultfontfeatures[\rmfamily]{Ligatures=TeX,Scale=1}
\fi
\usepackage{lmodern}
\ifPDFTeX\else
  % xetex/luatex font selection
\fi
% Use upquote if available, for straight quotes in verbatim environments
\IfFileExists{upquote.sty}{\usepackage{upquote}}{}
\IfFileExists{microtype.sty}{% use microtype if available
  \usepackage[]{microtype}
  \UseMicrotypeSet[protrusion]{basicmath} % disable protrusion for tt fonts
}{}
\makeatletter
\@ifundefined{KOMAClassName}{% if non-KOMA class
  \IfFileExists{parskip.sty}{%
    \usepackage{parskip}
  }{% else
    \setlength{\parindent}{0pt}
    \setlength{\parskip}{6pt plus 2pt minus 1pt}}
}{% if KOMA class
  \KOMAoptions{parskip=half}}
\makeatother
\usepackage{color}
\usepackage{fancyvrb}
\newcommand{\VerbBar}{|}
\newcommand{\VERB}{\Verb[commandchars=\\\{\}]}
\DefineVerbatimEnvironment{Highlighting}{Verbatim}{commandchars=\\\{\}}
% Add ',fontsize=\small' for more characters per line
\usepackage{framed}
\definecolor{shadecolor}{RGB}{248,248,248}
\newenvironment{Shaded}{\begin{snugshade}}{\end{snugshade}}
\newcommand{\AlertTok}[1]{\textcolor[rgb]{0.94,0.16,0.16}{#1}}
\newcommand{\AnnotationTok}[1]{\textcolor[rgb]{0.56,0.35,0.01}{\textbf{\textit{#1}}}}
\newcommand{\AttributeTok}[1]{\textcolor[rgb]{0.13,0.29,0.53}{#1}}
\newcommand{\BaseNTok}[1]{\textcolor[rgb]{0.00,0.00,0.81}{#1}}
\newcommand{\BuiltInTok}[1]{#1}
\newcommand{\CharTok}[1]{\textcolor[rgb]{0.31,0.60,0.02}{#1}}
\newcommand{\CommentTok}[1]{\textcolor[rgb]{0.56,0.35,0.01}{\textit{#1}}}
\newcommand{\CommentVarTok}[1]{\textcolor[rgb]{0.56,0.35,0.01}{\textbf{\textit{#1}}}}
\newcommand{\ConstantTok}[1]{\textcolor[rgb]{0.56,0.35,0.01}{#1}}
\newcommand{\ControlFlowTok}[1]{\textcolor[rgb]{0.13,0.29,0.53}{\textbf{#1}}}
\newcommand{\DataTypeTok}[1]{\textcolor[rgb]{0.13,0.29,0.53}{#1}}
\newcommand{\DecValTok}[1]{\textcolor[rgb]{0.00,0.00,0.81}{#1}}
\newcommand{\DocumentationTok}[1]{\textcolor[rgb]{0.56,0.35,0.01}{\textbf{\textit{#1}}}}
\newcommand{\ErrorTok}[1]{\textcolor[rgb]{0.64,0.00,0.00}{\textbf{#1}}}
\newcommand{\ExtensionTok}[1]{#1}
\newcommand{\FloatTok}[1]{\textcolor[rgb]{0.00,0.00,0.81}{#1}}
\newcommand{\FunctionTok}[1]{\textcolor[rgb]{0.13,0.29,0.53}{\textbf{#1}}}
\newcommand{\ImportTok}[1]{#1}
\newcommand{\InformationTok}[1]{\textcolor[rgb]{0.56,0.35,0.01}{\textbf{\textit{#1}}}}
\newcommand{\KeywordTok}[1]{\textcolor[rgb]{0.13,0.29,0.53}{\textbf{#1}}}
\newcommand{\NormalTok}[1]{#1}
\newcommand{\OperatorTok}[1]{\textcolor[rgb]{0.81,0.36,0.00}{\textbf{#1}}}
\newcommand{\OtherTok}[1]{\textcolor[rgb]{0.56,0.35,0.01}{#1}}
\newcommand{\PreprocessorTok}[1]{\textcolor[rgb]{0.56,0.35,0.01}{\textit{#1}}}
\newcommand{\RegionMarkerTok}[1]{#1}
\newcommand{\SpecialCharTok}[1]{\textcolor[rgb]{0.81,0.36,0.00}{\textbf{#1}}}
\newcommand{\SpecialStringTok}[1]{\textcolor[rgb]{0.31,0.60,0.02}{#1}}
\newcommand{\StringTok}[1]{\textcolor[rgb]{0.31,0.60,0.02}{#1}}
\newcommand{\VariableTok}[1]{\textcolor[rgb]{0.00,0.00,0.00}{#1}}
\newcommand{\VerbatimStringTok}[1]{\textcolor[rgb]{0.31,0.60,0.02}{#1}}
\newcommand{\WarningTok}[1]{\textcolor[rgb]{0.56,0.35,0.01}{\textbf{\textit{#1}}}}
\usepackage{graphicx}
\makeatletter
\newsavebox\pandoc@box
\newcommand*\pandocbounded[1]{% scales image to fit in text height/width
  \sbox\pandoc@box{#1}%
  \Gscale@div\@tempa{\textheight}{\dimexpr\ht\pandoc@box+\dp\pandoc@box\relax}%
  \Gscale@div\@tempb{\linewidth}{\wd\pandoc@box}%
  \ifdim\@tempb\p@<\@tempa\p@\let\@tempa\@tempb\fi% select the smaller of both
  \ifdim\@tempa\p@<\p@\scalebox{\@tempa}{\usebox\pandoc@box}%
  \else\usebox{\pandoc@box}%
  \fi%
}
% Set default figure placement to htbp
\def\fps@figure{htbp}
\makeatother
\setlength{\emergencystretch}{3em} % prevent overfull lines
\providecommand{\tightlist}{%
  \setlength{\itemsep}{0pt}\setlength{\parskip}{0pt}}
\usepackage{bookmark}
\IfFileExists{xurl.sty}{\usepackage{xurl}}{} % add URL line breaks if available
\urlstyle{same}
\hypersetup{
  pdftitle={CTA-ED Exercise 5: Unsupervised learning (topic models)},
  pdfauthor={Junke Guan},
  hidelinks,
  pdfcreator={LaTeX via pandoc}}

\title{CTA-ED Exercise 5: Unsupervised learning (topic models)}
\author{Junke Guan}
\date{11/03/2026}

\begin{document}
\maketitle

\subsection{Introduction}\label{introduction}

The hands-on exercise for this week focuses on: 1) estimating a topic
model ; 2) interpreting and visualizing results. Remember that you will
need to: 1) comment your code and 2) write out the interpretation of
your results.

You will learn how to:

\begin{itemize}
\tightlist
\item
  Generate document-term-matrices in format appropriate for topic
  modelling
\item
  Estimate a topic model using the \texttt{quanteda} and
  \texttt{topicmodels} package
\item
  Visualize results
\item
  Reverse engineer a test of model accuracy
\item
  Run some validation tests
\end{itemize}

\subsection{Setup}\label{setup}

Before proceeding, we'll load the packages we will need for this
tutorial.

\begin{Shaded}
\begin{Highlighting}[]
\FunctionTok{library}\NormalTok{(tidyverse) }\CommentTok{\# loads dplyr, ggplot2, and others}
\FunctionTok{library}\NormalTok{(stringr) }\CommentTok{\# to handle text elements}
\FunctionTok{library}\NormalTok{(tidytext) }\CommentTok{\# includes set of functions useful for manipulating text}
\FunctionTok{library}\NormalTok{(topicmodels) }\CommentTok{\# to estimate topic models}
\FunctionTok{library}\NormalTok{(gutenbergr) }\CommentTok{\# to get text data}
\FunctionTok{library}\NormalTok{(scales)}
\FunctionTok{library}\NormalTok{(tm)}
\FunctionTok{library}\NormalTok{(ggthemes) }\CommentTok{\# to make your plots look nice}
\FunctionTok{library}\NormalTok{(readr)}
\FunctionTok{library}\NormalTok{(quanteda)}
\FunctionTok{library}\NormalTok{(quanteda.textmodels)}
\end{Highlighting}
\end{Shaded}

You may need to install the preText package if you haven't done so yet.
For that you will need to run the next code chunk (it is currently set
to `eval=F', which tells R `do not execute this code chunk'). That
package is not readily available on through RStudio directly. It needs
to be downloaded from the Github repository set up by its creater
Matthew J Denny. We can do that using the command install\_github().
This command is part of the `devtools' package, which you will need to
install as well (if you haven't done so already). The devtools package
is directly available through R so it can be installed with the usual
command install\_packages.

\begin{Shaded}
\begin{Highlighting}[]
\CommentTok{\#install\_package(devtools)}
\NormalTok{devtools}\SpecialCharTok{::}\FunctionTok{install\_github}\NormalTok{(}\StringTok{"matthewjdenny/preText"}\NormalTok{)}
\FunctionTok{library}\NormalTok{(preText)}
\end{Highlighting}
\end{Shaded}

\section{Data collection}\label{data-collection}

We'll be using data from Alexis de Tocqueville's ``Democracy in
America.''

We have already downloaded some data for you, but we also included the
code to download it yourself (it is currently set to `eval=F' so it
won't run unless you remove the eval=F argument or you run the chunk
directly.

The code downloads these data, both Volume 1 and Volume 2, and combine
them into one data frame. For this, we'll be using the gutenbergr
package, which allows the user to download text data from over 60,000
out-of-copyright books. The ID for each book appears in the url for the
book selected after a search on \url{https://www.gutenberg.org/ebooks/}.

This example is adapted by \href{https://www.tidytextmining.com/}{Text
Mining with R: A Tidy Approach} by Julia Silge and David Robinson.

\includegraphics[width=1\linewidth,height=\textheight,keepaspectratio]{data/topicmodels/gutenberg.gif}

Here, we see that Volume of Tocqueville's ``Democracy in America'' is
stored as ``815''. A separate search reveals that Volume 2 is stored as
``816''.

\begin{Shaded}
\begin{Highlighting}[]
\NormalTok{tocq }\OtherTok{\textless{}{-}} \FunctionTok{gutenberg\_download}\NormalTok{(}\FunctionTok{c}\NormalTok{(}\DecValTok{815}\NormalTok{, }\DecValTok{816}\NormalTok{), }
                            \AttributeTok{meta\_fields =} \StringTok{"author"}\NormalTok{)}
\end{Highlighting}
\end{Shaded}

Or we can read the dataset we already downloaded for you in the
following way:

\begin{Shaded}
\begin{Highlighting}[]
\NormalTok{tocq  }\OtherTok{\textless{}{-}} \FunctionTok{readRDS}\NormalTok{(}\FunctionTok{gzcon}\NormalTok{(}\FunctionTok{url}\NormalTok{(}\StringTok{"https://github.com/cjbarrie/CTA{-}ED/blob/main/data/topicmodels/tocq.RDS?raw=true"}\NormalTok{)))}
\end{Highlighting}
\end{Shaded}

Once we have read in these data, we convert it into a different data
shape: the document-term-matrix. We also create a new columns, which we
call ``booknumber'' that recordss whether the term in question is from
Volume 1 or Volume 2. To convert from tidy into ``DocumentTermMatrix''
format we can first use \texttt{unnest\_tokens()} as we have done in
past exercises, remove stop words, and then use the \texttt{cast\_dtm()}
function to convert into a ``DocumentTermMatrix'' object.

\begin{Shaded}
\begin{Highlighting}[]
\NormalTok{tocq\_words }\OtherTok{\textless{}{-}}\NormalTok{ tocq }\SpecialCharTok{\%\textgreater{}\%}
  \FunctionTok{mutate}\NormalTok{(}\AttributeTok{booknumber =} \FunctionTok{ifelse}\NormalTok{(gutenberg\_id}\SpecialCharTok{==}\DecValTok{815}\NormalTok{, }\StringTok{"DiA1"}\NormalTok{, }\StringTok{"DiA2"}\NormalTok{)) }\SpecialCharTok{\%\textgreater{}\%}
  \FunctionTok{unnest\_tokens}\NormalTok{(word, text) }\SpecialCharTok{\%\textgreater{}\%}
  \FunctionTok{filter}\NormalTok{(}\SpecialCharTok{!}\FunctionTok{is.na}\NormalTok{(word)) }\SpecialCharTok{\%\textgreater{}\%}
  \FunctionTok{count}\NormalTok{(booknumber, word, }\AttributeTok{sort =} \ConstantTok{TRUE}\NormalTok{) }\SpecialCharTok{\%\textgreater{}\%}
  \FunctionTok{ungroup}\NormalTok{() }\SpecialCharTok{\%\textgreater{}\%}
  \FunctionTok{anti\_join}\NormalTok{(stop\_words)}
\end{Highlighting}
\end{Shaded}

\begin{verbatim}
## Joining with `by = join_by(word)`
\end{verbatim}

\begin{Shaded}
\begin{Highlighting}[]
\NormalTok{tocq\_dtm }\OtherTok{\textless{}{-}}\NormalTok{ tocq\_words }\SpecialCharTok{\%\textgreater{}\%}
  \FunctionTok{cast\_dtm}\NormalTok{(booknumber, word, n)}

\NormalTok{tm}\SpecialCharTok{::}\FunctionTok{inspect}\NormalTok{(tocq\_dtm)}
\end{Highlighting}
\end{Shaded}

\begin{verbatim}
## <<DocumentTermMatrix (documents: 2, terms: 12092)>>
## Non-/sparse entries: 17581/6603
## Sparsity           : 27%
## Maximal term length: 18
## Weighting          : term frequency (tf)
## Sample             :
##       Terms
## Docs   country democratic government laws nations people power society time
##   DiA1     357        212        556  397     233    516   543     290  311
##   DiA2     167        561        162  133     313    360   263     241  309
##       Terms
## Docs   united
##   DiA1    554
##   DiA2    227
\end{verbatim}

We see here that the data are now stored as a ``DocumentTermMatrix.'' In
this format, the matrix records the term (as equivalent of a column) and
the document (as equivalent of row), and the number of times the term
appears in the given document. Many terms will not appear in the
document, meaning that the matrix will be stored as ``sparse,'' meaning
there will be a preponderance of zeroes. Here, since we are looking only
at two documents that both come from a single volume set, the sparsity
is relatively low (only 27\%). In most applications, the sparsity will
be a lot higher, approaching 99\% or more.

Estimating our topic model is then relatively simple. All we need to do
if specify how many topics that we want to search for, and we can also
set our seed, which is needed to reproduce the same results each time
(as the model is a generative probabilistic one, meaning different
random iterations will produce different results).

\begin{Shaded}
\begin{Highlighting}[]
\NormalTok{tocq\_lda }\OtherTok{\textless{}{-}} \FunctionTok{LDA}\NormalTok{(tocq\_dtm, }\AttributeTok{k =} \DecValTok{10}\NormalTok{, }\AttributeTok{control =} \FunctionTok{list}\NormalTok{(}\AttributeTok{seed =} \DecValTok{1234}\NormalTok{))}
\end{Highlighting}
\end{Shaded}

After this we can extract the per-topic-per-word probabilities, called
``β'' from the model:

\begin{Shaded}
\begin{Highlighting}[]
\NormalTok{tocq\_topics }\OtherTok{\textless{}{-}} \FunctionTok{tidy}\NormalTok{(tocq\_lda, }\AttributeTok{matrix =} \StringTok{"beta"}\NormalTok{)}

\FunctionTok{head}\NormalTok{(tocq\_topics, }\AttributeTok{n =} \DecValTok{10}\NormalTok{)}
\end{Highlighting}
\end{Shaded}

\begin{verbatim}
## # A tibble: 10 x 3
##    topic term          beta
##    <int> <chr>        <dbl>
##  1     1 democratic 0.00855
##  2     2 democratic 0.0115 
##  3     3 democratic 0.00444
##  4     4 democratic 0.0193 
##  5     5 democratic 0.00254
##  6     6 democratic 0.00866
##  7     7 democratic 0.00165
##  8     8 democratic 0.0108 
##  9     9 democratic 0.00276
## 10    10 democratic 0.00334
\end{verbatim}

We now have data stored as one topic-per-term-per-row. The betas listed
here represent the probability that the given term belongs to a given
topic. So, here, we see that the term ``democratic'' is most likely to
belong to topic 4. Strictly, this probability represents the probability
that the term is generated from the topic in question.

We can then plots the top terms, in terms of beta, for each topic as
follows:

\begin{Shaded}
\begin{Highlighting}[]
\NormalTok{tocq\_top\_terms }\OtherTok{\textless{}{-}}\NormalTok{ tocq\_topics }\SpecialCharTok{\%\textgreater{}\%}
  \FunctionTok{group\_by}\NormalTok{(topic) }\SpecialCharTok{\%\textgreater{}\%}
  \FunctionTok{top\_n}\NormalTok{(}\DecValTok{10}\NormalTok{, beta) }\SpecialCharTok{\%\textgreater{}\%}
  \FunctionTok{ungroup}\NormalTok{() }\SpecialCharTok{\%\textgreater{}\%}
  \FunctionTok{arrange}\NormalTok{(topic, }\SpecialCharTok{{-}}\NormalTok{beta)}

\NormalTok{tocq\_top\_terms }\SpecialCharTok{\%\textgreater{}\%}
  \FunctionTok{mutate}\NormalTok{(}\AttributeTok{term =} \FunctionTok{reorder\_within}\NormalTok{(term, beta, topic)) }\SpecialCharTok{\%\textgreater{}\%}
  \FunctionTok{ggplot}\NormalTok{(}\FunctionTok{aes}\NormalTok{(beta, term, }\AttributeTok{fill =} \FunctionTok{factor}\NormalTok{(topic))) }\SpecialCharTok{+}
  \FunctionTok{geom\_col}\NormalTok{(}\AttributeTok{show.legend =} \ConstantTok{FALSE}\NormalTok{) }\SpecialCharTok{+}
  \FunctionTok{facet\_wrap}\NormalTok{(}\SpecialCharTok{\textasciitilde{}}\NormalTok{ topic, }\AttributeTok{scales =} \StringTok{"free"}\NormalTok{, }\AttributeTok{ncol =} \DecValTok{4}\NormalTok{) }\SpecialCharTok{+}
  \FunctionTok{scale\_y\_reordered}\NormalTok{() }\SpecialCharTok{+}
  \FunctionTok{theme\_tufte}\NormalTok{(}\AttributeTok{base\_family =} \StringTok{"Helvetica"}\NormalTok{)}
\end{Highlighting}
\end{Shaded}

\pandocbounded{\includegraphics[keepaspectratio]{unsupervised-topicmodels_files/figure-latex/unnamed-chunk-8-1.pdf}}

But how do we actually evaluate these topics? Here, the topics all seem
pretty similar.

\subsection{Evaluating topic model}\label{evaluating-topic-model}

Well, one way to evaluate the performance of unspervised forms of
classification is by testing our model on an outcome that is already
known.

Here, two topics that are most obvious are the `topics' Volume 1 and
Volume 2 of Tocqueville's ``Democracy in America.'' Volume 1 of
Tocqueville's work deals more obviously with abstract constitutional
ideas and questions of race; Volume 2 focuses on more esoteric aspects
of American society. Listen an ``In Our Time'' episode with Melvyn Bragg
discussing Democracy in America
\href{https://www.bbc.co.uk/programmes/b09vyw0x}{here}.

Given these differences in focus, we might think that a generative model
could accurately assign to topic (i.e., Volume) with some accuracy.

\subsubsection{Plot relative word
frequencies}\label{plot-relative-word-frequencies}

First let's have a look and see whether there really are words obviously
distinguishing the two Volumes.

\begin{Shaded}
\begin{Highlighting}[]
\NormalTok{tidy\_tocq }\OtherTok{\textless{}{-}}\NormalTok{ tocq }\SpecialCharTok{\%\textgreater{}\%}
  \FunctionTok{unnest\_tokens}\NormalTok{(word, text) }\SpecialCharTok{\%\textgreater{}\%}
  \FunctionTok{anti\_join}\NormalTok{(stop\_words)}
\end{Highlighting}
\end{Shaded}

\begin{verbatim}
## Joining with `by = join_by(word)`
\end{verbatim}

\begin{Shaded}
\begin{Highlighting}[]
\DocumentationTok{\#\# Count most common words in both}
\NormalTok{tidy\_tocq }\SpecialCharTok{\%\textgreater{}\%}
  \FunctionTok{count}\NormalTok{(word, }\AttributeTok{sort =} \ConstantTok{TRUE}\NormalTok{)}
\end{Highlighting}
\end{Shaded}

\begin{verbatim}
## # A tibble: 12,092 x 2
##    word           n
##    <chr>      <int>
##  1 people       876
##  2 power        806
##  3 united       781
##  4 democratic   773
##  5 government   718
##  6 time         620
##  7 nations      546
##  8 society      531
##  9 laws         530
## 10 country      524
## # i 12,082 more rows
\end{verbatim}

\begin{Shaded}
\begin{Highlighting}[]
\NormalTok{bookfreq }\OtherTok{\textless{}{-}}\NormalTok{ tidy\_tocq }\SpecialCharTok{\%\textgreater{}\%}
  \FunctionTok{mutate}\NormalTok{(}\AttributeTok{booknumber =} \FunctionTok{ifelse}\NormalTok{(gutenberg\_id}\SpecialCharTok{==}\DecValTok{815}\NormalTok{, }\StringTok{"DiA1"}\NormalTok{, }\StringTok{"DiA2"}\NormalTok{)) }\SpecialCharTok{\%\textgreater{}\%}
  \FunctionTok{mutate}\NormalTok{(}\AttributeTok{word =} \FunctionTok{str\_extract}\NormalTok{(word, }\StringTok{"[a{-}z\textquotesingle{}]+"}\NormalTok{)) }\SpecialCharTok{\%\textgreater{}\%}
  \FunctionTok{count}\NormalTok{(booknumber, word) }\SpecialCharTok{\%\textgreater{}\%}
  \FunctionTok{group\_by}\NormalTok{(booknumber) }\SpecialCharTok{\%\textgreater{}\%}
  \FunctionTok{mutate}\NormalTok{(}\AttributeTok{proportion =}\NormalTok{ n }\SpecialCharTok{/} \FunctionTok{sum}\NormalTok{(n)) }\SpecialCharTok{\%\textgreater{}\%} 
  \FunctionTok{select}\NormalTok{(}\SpecialCharTok{{-}}\NormalTok{n) }\SpecialCharTok{\%\textgreater{}\%} 
  \FunctionTok{spread}\NormalTok{(booknumber, proportion)}

\FunctionTok{ggplot}\NormalTok{(bookfreq, }\FunctionTok{aes}\NormalTok{(}\AttributeTok{x =}\NormalTok{ DiA1, }\AttributeTok{y =}\NormalTok{ DiA2, }\AttributeTok{color =} \FunctionTok{abs}\NormalTok{(DiA1 }\SpecialCharTok{{-}}\NormalTok{ DiA2))) }\SpecialCharTok{+}
  \FunctionTok{geom\_abline}\NormalTok{(}\AttributeTok{color =} \StringTok{"gray40"}\NormalTok{, }\AttributeTok{lty =} \DecValTok{2}\NormalTok{) }\SpecialCharTok{+}
  \FunctionTok{geom\_jitter}\NormalTok{(}\AttributeTok{alpha =} \FloatTok{0.1}\NormalTok{, }\AttributeTok{size =} \FloatTok{2.5}\NormalTok{, }\AttributeTok{width =} \FloatTok{0.3}\NormalTok{, }\AttributeTok{height =} \FloatTok{0.3}\NormalTok{) }\SpecialCharTok{+}
  \FunctionTok{geom\_text}\NormalTok{(}\FunctionTok{aes}\NormalTok{(}\AttributeTok{label =}\NormalTok{ word), }\AttributeTok{check\_overlap =} \ConstantTok{TRUE}\NormalTok{, }\AttributeTok{vjust =} \FloatTok{1.5}\NormalTok{) }\SpecialCharTok{+}
  \FunctionTok{scale\_x\_log10}\NormalTok{(}\AttributeTok{labels =} \FunctionTok{percent\_format}\NormalTok{()) }\SpecialCharTok{+}
  \FunctionTok{scale\_y\_log10}\NormalTok{(}\AttributeTok{labels =} \FunctionTok{percent\_format}\NormalTok{()) }\SpecialCharTok{+}
  \FunctionTok{scale\_color\_gradient}\NormalTok{(}\AttributeTok{limits =} \FunctionTok{c}\NormalTok{(}\DecValTok{0}\NormalTok{, }\FloatTok{0.001}\NormalTok{), }\AttributeTok{low =} \StringTok{"darkslategray4"}\NormalTok{, }\AttributeTok{high =} \StringTok{"gray75"}\NormalTok{) }\SpecialCharTok{+}
  \FunctionTok{theme\_tufte}\NormalTok{(}\AttributeTok{base\_family =} \StringTok{"Helvetica"}\NormalTok{) }\SpecialCharTok{+}
  \FunctionTok{theme}\NormalTok{(}\AttributeTok{legend.position=}\StringTok{"none"}\NormalTok{, }
        \AttributeTok{strip.background =} \FunctionTok{element\_blank}\NormalTok{(), }
        \AttributeTok{strip.text.x =} \FunctionTok{element\_blank}\NormalTok{()) }\SpecialCharTok{+}
  \FunctionTok{labs}\NormalTok{(}\AttributeTok{x =} \StringTok{"Tocqueville DiA 2"}\NormalTok{, }\AttributeTok{y =} \StringTok{"Tocqueville DiA 1"}\NormalTok{) }\SpecialCharTok{+}
  \FunctionTok{coord\_equal}\NormalTok{()}
\end{Highlighting}
\end{Shaded}

\begin{verbatim}
## Warning: Removed 6173 rows containing missing values or values outside the scale range
## (`geom_point()`).
\end{verbatim}

\begin{verbatim}
## Warning: Removed 6174 rows containing missing values or values outside the scale range
## (`geom_text()`).
\end{verbatim}

\pandocbounded{\includegraphics[keepaspectratio]{unsupervised-topicmodels_files/figure-latex/unnamed-chunk-9-1.pdf}}

We see that there do seem to be some marked distinguishing
characteristics. In the plot above, for example, we see that more
abstract notions of state systems appear with greater frequency in
Volume 1 while Volume 2 seems to contain words specific to America
(e.g., ``north'' and ``south'') with greater frequency. The way to read
the above plot is that words positioned further away from the diagonal
line appear with greater frequency in one volume versus the other.

\subsubsection{Split into chapter
documents}\label{split-into-chapter-documents}

In the below, we first separate the volumes into chapters, then we
repeat the same procedure as above. The only difference now is that
instead of two documents representing the two full volumes of
Tocqueville's work, we now have 132 documents, each representing an
individual chapter. Notice now that the sparsity is much increased:
around 96\%.

\begin{Shaded}
\begin{Highlighting}[]
\NormalTok{tocq }\OtherTok{\textless{}{-}}\NormalTok{ tocq }\SpecialCharTok{\%\textgreater{}\%}
  \FunctionTok{filter}\NormalTok{(}\SpecialCharTok{!}\FunctionTok{is.na}\NormalTok{(text))}

\CommentTok{\# Divide into documents, each representing one chapter}
\NormalTok{tocq\_chapter }\OtherTok{\textless{}{-}}\NormalTok{ tocq }\SpecialCharTok{\%\textgreater{}\%}
  \FunctionTok{mutate}\NormalTok{(}\AttributeTok{booknumber =} \FunctionTok{ifelse}\NormalTok{(gutenberg\_id}\SpecialCharTok{==}\DecValTok{815}\NormalTok{, }\StringTok{"DiA1"}\NormalTok{, }\StringTok{"DiA2"}\NormalTok{)) }\SpecialCharTok{\%\textgreater{}\%}
  \FunctionTok{group\_by}\NormalTok{(booknumber) }\SpecialCharTok{\%\textgreater{}\%}
  \FunctionTok{mutate}\NormalTok{(}\AttributeTok{chapter =} \FunctionTok{cumsum}\NormalTok{(}\FunctionTok{str\_detect}\NormalTok{(text, }\FunctionTok{regex}\NormalTok{(}\StringTok{"\^{}chapter "}\NormalTok{, }\AttributeTok{ignore\_case =} \ConstantTok{TRUE}\NormalTok{)))) }\SpecialCharTok{\%\textgreater{}\%}
  \FunctionTok{ungroup}\NormalTok{() }\SpecialCharTok{\%\textgreater{}\%}
  \FunctionTok{filter}\NormalTok{(chapter }\SpecialCharTok{\textgreater{}} \DecValTok{0}\NormalTok{) }\SpecialCharTok{\%\textgreater{}\%}
  \FunctionTok{unite}\NormalTok{(document, booknumber, chapter)}

\CommentTok{\# Split into words}
\NormalTok{tocq\_chapter\_word }\OtherTok{\textless{}{-}}\NormalTok{ tocq\_chapter }\SpecialCharTok{\%\textgreater{}\%}
  \FunctionTok{unnest\_tokens}\NormalTok{(word, text)}

\CommentTok{\# Find document{-}word counts}
\NormalTok{tocq\_word\_counts }\OtherTok{\textless{}{-}}\NormalTok{ tocq\_chapter\_word }\SpecialCharTok{\%\textgreater{}\%}
  \FunctionTok{anti\_join}\NormalTok{(stop\_words) }\SpecialCharTok{\%\textgreater{}\%}
  \FunctionTok{count}\NormalTok{(document, word, }\AttributeTok{sort =} \ConstantTok{TRUE}\NormalTok{) }\SpecialCharTok{\%\textgreater{}\%}
  \FunctionTok{ungroup}\NormalTok{()}
\end{Highlighting}
\end{Shaded}

\begin{verbatim}
## Joining with `by = join_by(word)`
\end{verbatim}

\begin{Shaded}
\begin{Highlighting}[]
\NormalTok{tocq\_word\_counts}
\end{Highlighting}
\end{Shaded}

\begin{verbatim}
## # A tibble: 69,781 x 3
##    document word             n
##    <chr>    <chr>        <int>
##  1 DiA2_76  united          88
##  2 DiA2_60  honor           70
##  3 DiA1_52  union           66
##  4 DiA2_76  president       60
##  5 DiA2_76  law             59
##  6 DiA1_42  jury            57
##  7 DiA2_76  time            50
##  8 DiA1_11  township        49
##  9 DiA1_21  federal         48
## 10 DiA2_76  constitution    48
## # i 69,771 more rows
\end{verbatim}

\begin{Shaded}
\begin{Highlighting}[]
\CommentTok{\# Cast into DTM format for LDA analysis}

\NormalTok{tocq\_chapters\_dtm }\OtherTok{\textless{}{-}}\NormalTok{ tocq\_word\_counts }\SpecialCharTok{\%\textgreater{}\%}
  \FunctionTok{cast\_dtm}\NormalTok{(document, word, n)}

\NormalTok{tm}\SpecialCharTok{::}\FunctionTok{inspect}\NormalTok{(tocq\_chapters\_dtm)}
\end{Highlighting}
\end{Shaded}

\begin{verbatim}
## <<DocumentTermMatrix (documents: 132, terms: 11898)>>
## Non-/sparse entries: 69781/1500755
## Sparsity           : 96%
## Maximal term length: 18
## Weighting          : term frequency (tf)
## Sample             :
##          Terms
## Docs      country democratic government laws nations people power public time
##   DiA1_11      10          0         23   19       7     13    19     15    6
##   DiA1_13      13          5         34    9      12     17    37     15    6
##   DiA1_20       9          0         25   13       2     14    32     13   10
##   DiA1_21       4          0         20   29       6     12    20      5    5
##   DiA1_23      10          0         35    9      24     20    13      4    8
##   DiA1_31       7         12         10   13       4     30    18     31    6
##   DiA1_32      10         14         25    6       9     25    11     43    8
##   DiA1_47      12          2          5    3       3      6     8      0    3
##   DiA1_56      12          0          3    7      19      3     8      3   22
##   DiA2_76      11         10         24   39      12     31    27     27   50
##          Terms
## Docs      united
##   DiA1_11     13
##   DiA1_13     19
##   DiA1_20     21
##   DiA1_21     23
##   DiA1_23     15
##   DiA1_31     11
##   DiA1_32     14
##   DiA1_47      8
##   DiA1_56     25
##   DiA2_76     88
\end{verbatim}

We then re-estimate the topic model with this new DocumentTermMatrix
object, specifying k equal to 2. This will enable us to evaluate whether
a topic model is able to generatively assign to volume with accuracy.

\begin{Shaded}
\begin{Highlighting}[]
\NormalTok{tocq\_chapters\_lda }\OtherTok{\textless{}{-}} \FunctionTok{LDA}\NormalTok{(tocq\_chapters\_dtm, }\AttributeTok{k =} \DecValTok{2}\NormalTok{, }\AttributeTok{control =} \FunctionTok{list}\NormalTok{(}\AttributeTok{seed =} \DecValTok{1234}\NormalTok{))}
\end{Highlighting}
\end{Shaded}

After this, it is worth looking at another output of the latent
dirichlet allocation procedure. The γ probability represents the
per-document-per-topic probability or, in other words, the probability
that a given document (here: chapter) belongs to a particular topic (and
here, we are assuming these topics represent volumes).

The gamma values are therefore the estimated proportion of words within
a given chapter allocated to a given volume.

\begin{Shaded}
\begin{Highlighting}[]
\NormalTok{tocq\_chapters\_gamma }\OtherTok{\textless{}{-}} \FunctionTok{tidy}\NormalTok{(tocq\_chapters\_lda, }\AttributeTok{matrix =} \StringTok{"gamma"}\NormalTok{)}
\NormalTok{tocq\_chapters\_gamma}
\end{Highlighting}
\end{Shaded}

\begin{verbatim}
## # A tibble: 264 x 3
##    document topic     gamma
##    <chr>    <int>     <dbl>
##  1 DiA2_76      1 0.551    
##  2 DiA2_60      1 1.000    
##  3 DiA1_52      1 0.0000464
##  4 DiA1_42      1 0.0000746
##  5 DiA1_11      1 0.0000382
##  6 DiA1_21      1 0.0000437
##  7 DiA1_20      1 0.0000425
##  8 DiA1_28      1 0.249    
##  9 DiA1_50      1 0.0000477
## 10 DiA1_22      1 0.0000466
## # i 254 more rows
\end{verbatim}

\subsubsection{Examine consensus}\label{examine-consensus}

Now that we have these topic probabilities, we can see how well our
unsupervised learning did at distinguishing the two volumes generatively
just from the words contained in each chapter.

\begin{Shaded}
\begin{Highlighting}[]
\CommentTok{\# First separate the document name into title and chapter}

\NormalTok{tocq\_chapters\_gamma }\OtherTok{\textless{}{-}}\NormalTok{ tocq\_chapters\_gamma }\SpecialCharTok{\%\textgreater{}\%}
  \FunctionTok{separate}\NormalTok{(document, }\FunctionTok{c}\NormalTok{(}\StringTok{"title"}\NormalTok{, }\StringTok{"chapter"}\NormalTok{), }\AttributeTok{sep =} \StringTok{"\_"}\NormalTok{, }\AttributeTok{convert =} \ConstantTok{TRUE}\NormalTok{)}

\NormalTok{tocq\_chapter\_classifications }\OtherTok{\textless{}{-}}\NormalTok{ tocq\_chapters\_gamma }\SpecialCharTok{\%\textgreater{}\%}
  \FunctionTok{group\_by}\NormalTok{(title, chapter) }\SpecialCharTok{\%\textgreater{}\%}
  \FunctionTok{top\_n}\NormalTok{(}\DecValTok{1}\NormalTok{, gamma) }\SpecialCharTok{\%\textgreater{}\%}
  \FunctionTok{ungroup}\NormalTok{()}

\NormalTok{tocq\_book\_topics }\OtherTok{\textless{}{-}}\NormalTok{ tocq\_chapter\_classifications }\SpecialCharTok{\%\textgreater{}\%}
  \FunctionTok{count}\NormalTok{(title, topic) }\SpecialCharTok{\%\textgreater{}\%}
  \FunctionTok{group\_by}\NormalTok{(title) }\SpecialCharTok{\%\textgreater{}\%}
  \FunctionTok{top\_n}\NormalTok{(}\DecValTok{1}\NormalTok{, n) }\SpecialCharTok{\%\textgreater{}\%}
  \FunctionTok{ungroup}\NormalTok{() }\SpecialCharTok{\%\textgreater{}\%}
  \FunctionTok{transmute}\NormalTok{(}\AttributeTok{consensus =}\NormalTok{ title, topic)}

\NormalTok{tocq\_chapter\_classifications }\SpecialCharTok{\%\textgreater{}\%}
  \FunctionTok{inner\_join}\NormalTok{(tocq\_book\_topics, }\AttributeTok{by =} \StringTok{"topic"}\NormalTok{) }\SpecialCharTok{\%\textgreater{}\%}
  \FunctionTok{filter}\NormalTok{(title }\SpecialCharTok{!=}\NormalTok{ consensus)}
\end{Highlighting}
\end{Shaded}

\begin{verbatim}
## # A tibble: 15 x 5
##    title chapter topic gamma consensus
##    <chr>   <int> <int> <dbl> <chr>    
##  1 DiA1       45     1 0.762 DiA2     
##  2 DiA1        5     1 0.504 DiA2     
##  3 DiA1       33     1 0.570 DiA2     
##  4 DiA1       34     1 0.626 DiA2     
##  5 DiA1       41     1 0.512 DiA2     
##  6 DiA1       44     1 0.765 DiA2     
##  7 DiA1        8     1 0.791 DiA2     
##  8 DiA1        4     1 0.717 DiA2     
##  9 DiA1       35     1 0.576 DiA2     
## 10 DiA1       39     1 0.577 DiA2     
## 11 DiA1        7     1 0.687 DiA2     
## 12 DiA1       29     1 0.983 DiA2     
## 13 DiA1        6     1 0.707 DiA2     
## 14 DiA2       27     2 0.654 DiA1     
## 15 DiA2       21     2 0.510 DiA1
\end{verbatim}

\begin{Shaded}
\begin{Highlighting}[]
\CommentTok{\# Look document{-}word pairs were to see which words in each documents were assigned}
\CommentTok{\# to a given topic}

\NormalTok{assignments }\OtherTok{\textless{}{-}} \FunctionTok{augment}\NormalTok{(tocq\_chapters\_lda, }\AttributeTok{data =}\NormalTok{ tocq\_chapters\_dtm)}
\NormalTok{assignments}
\end{Highlighting}
\end{Shaded}

\begin{verbatim}
## # A tibble: 69,781 x 4
##    document term   count .topic
##    <chr>    <chr>  <dbl>  <dbl>
##  1 DiA2_76  united    88      2
##  2 DiA2_60  united     6      1
##  3 DiA1_52  united    11      2
##  4 DiA1_42  united     7      2
##  5 DiA1_11  united    13      2
##  6 DiA1_21  united    23      2
##  7 DiA1_20  united    21      2
##  8 DiA1_28  united    14      2
##  9 DiA1_50  united     5      2
## 10 DiA1_22  united     8      2
## # i 69,771 more rows
\end{verbatim}

\begin{Shaded}
\begin{Highlighting}[]
\NormalTok{assignments }\OtherTok{\textless{}{-}}\NormalTok{ assignments }\SpecialCharTok{\%\textgreater{}\%}
  \FunctionTok{separate}\NormalTok{(document, }\FunctionTok{c}\NormalTok{(}\StringTok{"title"}\NormalTok{, }\StringTok{"chapter"}\NormalTok{), }\AttributeTok{sep =} \StringTok{"\_"}\NormalTok{, }\AttributeTok{convert =} \ConstantTok{TRUE}\NormalTok{) }\SpecialCharTok{\%\textgreater{}\%}
  \FunctionTok{inner\_join}\NormalTok{(tocq\_book\_topics, }\AttributeTok{by =} \FunctionTok{c}\NormalTok{(}\StringTok{".topic"} \OtherTok{=} \StringTok{"topic"}\NormalTok{))}

\NormalTok{assignments }\SpecialCharTok{\%\textgreater{}\%}
  \FunctionTok{count}\NormalTok{(title, consensus, }\AttributeTok{wt =}\NormalTok{ count) }\SpecialCharTok{\%\textgreater{}\%}
  \FunctionTok{group\_by}\NormalTok{(title) }\SpecialCharTok{\%\textgreater{}\%}
  \FunctionTok{mutate}\NormalTok{(}\AttributeTok{percent =}\NormalTok{ n }\SpecialCharTok{/} \FunctionTok{sum}\NormalTok{(n)) }\SpecialCharTok{\%\textgreater{}\%}
  \FunctionTok{ggplot}\NormalTok{(}\FunctionTok{aes}\NormalTok{(consensus, title, }\AttributeTok{fill =}\NormalTok{ percent)) }\SpecialCharTok{+}
  \FunctionTok{geom\_tile}\NormalTok{() }\SpecialCharTok{+}
  \FunctionTok{scale\_fill\_gradient2}\NormalTok{(}\AttributeTok{high =} \StringTok{"red"}\NormalTok{, }\AttributeTok{label =} \FunctionTok{percent\_format}\NormalTok{()) }\SpecialCharTok{+}
  \FunctionTok{geom\_text}\NormalTok{(}\FunctionTok{aes}\NormalTok{(}\AttributeTok{x =}\NormalTok{ consensus, }\AttributeTok{y =}\NormalTok{ title, }\AttributeTok{label =}\NormalTok{ scales}\SpecialCharTok{::}\FunctionTok{percent}\NormalTok{(percent))) }\SpecialCharTok{+}
  \FunctionTok{theme\_tufte}\NormalTok{(}\AttributeTok{base\_family =} \StringTok{"Helvetica"}\NormalTok{) }\SpecialCharTok{+}
  \FunctionTok{theme}\NormalTok{(}\AttributeTok{axis.text.x =} \FunctionTok{element\_text}\NormalTok{(}\AttributeTok{angle =} \DecValTok{90}\NormalTok{, }\AttributeTok{hjust =} \DecValTok{1}\NormalTok{),}
        \AttributeTok{panel.grid =} \FunctionTok{element\_blank}\NormalTok{()) }\SpecialCharTok{+}
  \FunctionTok{labs}\NormalTok{(}\AttributeTok{x =} \StringTok{"Book words assigned to"}\NormalTok{,}
       \AttributeTok{y =} \StringTok{"Book words came from"}\NormalTok{,}
       \AttributeTok{fill =} \StringTok{"\% of assignments"}\NormalTok{)}
\end{Highlighting}
\end{Shaded}

\pandocbounded{\includegraphics[keepaspectratio]{unsupervised-topicmodels_files/figure-latex/unnamed-chunk-13-1.pdf}}

Not bad! We see that the model estimated with accuracy 91\% of chapters
in Volume 2 and 79\% of chapters in Volume 1

\subsection{Validation}\label{validation}

In the articles by @ying\_topics\_2021 and @denny\_text\_2018 from this
and previous weeks, we read about potential validation techniques.

In this section, we'll be using the \texttt{preText} package mentioned
in @denny\_text\_2018 to see the impact of different pre-processing
choices on our text. Here, I am adapting from a
\href{http://www.mjdenny.com/getting_started_with_preText.html}{tutorial}
by Matthew Denny.

First we need to reformat our text into a \texttt{quanteda} corpus
object.

\begin{Shaded}
\begin{Highlighting}[]
\CommentTok{\# load in corpus of Tocequeville text data.}
\NormalTok{corp }\OtherTok{\textless{}{-}} \FunctionTok{corpus}\NormalTok{(tocq, }\AttributeTok{text\_field =} \StringTok{"text"}\NormalTok{)}
\CommentTok{\# use first 10 documents for example}
\NormalTok{documents }\OtherTok{\textless{}{-}}\NormalTok{ corp[}\FunctionTok{sample}\NormalTok{(}\DecValTok{1}\SpecialCharTok{:}\DecValTok{30000}\NormalTok{,}\DecValTok{1000}\NormalTok{)]}
\CommentTok{\# take a look at the document names}
\FunctionTok{print}\NormalTok{(}\FunctionTok{names}\NormalTok{(documents[}\DecValTok{1}\SpecialCharTok{:}\DecValTok{10}\NormalTok{]))}
\end{Highlighting}
\end{Shaded}

\begin{verbatim}
##  [1] "text9079"  "text4245"  "text8249"  "text16921" "text11970" "text21533"
##  [7] "text3033"  "text4502"  "text17293" "text26455"
\end{verbatim}

And now we are ready to preprocess in different ways. Here, we are
including n-grams so we are preprocessing the text in 128 different
ways. This takes about ten minutes to run on a machine with 8GB RAM.

\begin{Shaded}
\begin{Highlighting}[]
\NormalTok{preprocessed\_documents }\OtherTok{\textless{}{-}} \FunctionTok{factorial\_preprocessing}\NormalTok{(}
\NormalTok{    documents,}
    \AttributeTok{use\_ngrams =} \ConstantTok{TRUE}\NormalTok{,}
    \AttributeTok{infrequent\_term\_threshold =} \FloatTok{0.2}\NormalTok{,}
    \AttributeTok{verbose =} \ConstantTok{FALSE}\NormalTok{)}
\end{Highlighting}
\end{Shaded}

We can then get the results of our pre-processing, comparing the
distance between documents that have been processed in different ways.

\begin{Shaded}
\begin{Highlighting}[]
\NormalTok{preText\_results }\OtherTok{\textless{}{-}} \FunctionTok{preText}\NormalTok{(}
\NormalTok{    preprocessed\_documents,}
    \AttributeTok{dataset\_name =} \StringTok{"Tocqueville text"}\NormalTok{,}
    \AttributeTok{distance\_method =} \StringTok{"cosine"}\NormalTok{,}
    \AttributeTok{num\_comparisons =} \DecValTok{20}\NormalTok{,}
    \AttributeTok{verbose =} \ConstantTok{FALSE}\NormalTok{)}
\end{Highlighting}
\end{Shaded}

And we can plot these accordingly.

\begin{Shaded}
\begin{Highlighting}[]
\FunctionTok{preText\_score\_plot}\NormalTok{(preText\_results)}
\end{Highlighting}
\end{Shaded}

\includegraphics[width=1\linewidth,height=\textheight,keepaspectratio]{data/topicmodels/pretext_results.png}

\subsection{Exercises}\label{exercises}

\begin{enumerate}
\def\labelenumi{\arabic{enumi}.}
\tightlist
\item
  Choose another book or set of books from Project Gutenberg
\item
  Run your own topic model on these books, changing the k of topics, and
  evaluating accuracy.
\item
  Validate different pre-processing techniques using \texttt{preText} on
  the new book(s) of your choice.
\end{enumerate}

For this exercise, I analyse Jane Austen's \emph{Pride and Prejudice}
(Project Gutenberg ID: 1342). The novel is divided into chapters, which
are treated as individual documents for topic modelling.

\begin{Shaded}
\begin{Highlighting}[]
\NormalTok{book }\OtherTok{\textless{}{-}} \FunctionTok{gutenberg\_download}\NormalTok{(}\DecValTok{1342}\NormalTok{, }\AttributeTok{meta\_fields =} \StringTok{"author"}\NormalTok{)}
\end{Highlighting}
\end{Shaded}

\begin{verbatim}
## Using mirror https://aleph.pglaf.org.
\end{verbatim}

\begin{Shaded}
\begin{Highlighting}[]
\FunctionTok{head}\NormalTok{(book)}
\end{Highlighting}
\end{Shaded}

\begin{verbatim}
## # A tibble: 6 x 3
##   gutenberg_id text                                             author      
##          <int> <chr>                                            <chr>       
## 1         1342 "                            [Illustration:"     Austen, Jane
## 2         1342 ""                                               Austen, Jane
## 3         1342 "                             GEORGE ALLEN"      Austen, Jane
## 4         1342 "                               PUBLISHER"       Austen, Jane
## 5         1342 ""                                               Austen, Jane
## 6         1342 "                        156 CHARING CROSS ROAD" Austen, Jane
\end{verbatim}

\begin{Shaded}
\begin{Highlighting}[]
\NormalTok{book\_chapters }\OtherTok{\textless{}{-}}\NormalTok{ book }\SpecialCharTok{\%\textgreater{}\%}
  \FunctionTok{filter}\NormalTok{(}\SpecialCharTok{!}\FunctionTok{is.na}\NormalTok{(text)) }\SpecialCharTok{\%\textgreater{}\%}
  \FunctionTok{mutate}\NormalTok{(}\AttributeTok{chapter =} \FunctionTok{cumsum}\NormalTok{(}\FunctionTok{str\_detect}\NormalTok{(text, }\FunctionTok{regex}\NormalTok{(}\StringTok{"\^{}chapter"}\NormalTok{, }\AttributeTok{ignore\_case =} \ConstantTok{TRUE}\NormalTok{)))) }\SpecialCharTok{\%\textgreater{}\%}
  \FunctionTok{filter}\NormalTok{(chapter }\SpecialCharTok{\textgreater{}} \DecValTok{0}\NormalTok{)}

\FunctionTok{head}\NormalTok{(book\_chapters)}
\end{Highlighting}
\end{Shaded}

\begin{verbatim}
## # A tibble: 6 x 4
##   gutenberg_id text                                               author chapter
##          <int> <chr>                                              <chr>    <int>
## 1         1342 "Chapter I.]"                                      Auste~       1
## 2         1342 ""                                                 Auste~       1
## 3         1342 ""                                                 Auste~       1
## 4         1342 "It is a truth universally acknowledged, that a s~ Auste~       1
## 5         1342 "of a good fortune must be in want of a wife."     Auste~       1
## 6         1342 ""                                                 Auste~       1
\end{verbatim}

To estimate a topic model, the text needs to be divided into multiple
documents. In this analysis, each chapter of \emph{Pride and Prejudice}
is treated as a separate document. This approach allows the model to
identify thematic differences across chapters of the novel.

\begin{Shaded}
\begin{Highlighting}[]
\NormalTok{book\_words }\OtherTok{\textless{}{-}}\NormalTok{ book\_chapters }\SpecialCharTok{\%\textgreater{}\%}
  \FunctionTok{unnest\_tokens}\NormalTok{(word, text)}

\FunctionTok{head}\NormalTok{(book\_words)}
\end{Highlighting}
\end{Shaded}

\begin{verbatim}
## # A tibble: 6 x 4
##   gutenberg_id author       chapter word   
##          <int> <chr>          <int> <chr>  
## 1         1342 Austen, Jane       1 chapter
## 2         1342 Austen, Jane       1 i      
## 3         1342 Austen, Jane       1 it     
## 4         1342 Austen, Jane       1 is     
## 5         1342 Austen, Jane       1 a      
## 6         1342 Austen, Jane       1 truth
\end{verbatim}

Next, the text is tokenized into individual words using the
\texttt{unnest\_tokens()} function from the tidytext package. This
converts the text into a tidy format where each row represents a single
word. Tokenization is an important preprocessing step in text analysis,
as it allows us to analyse word frequencies and construct a
document-term matrix for topic modelling.

\begin{Shaded}
\begin{Highlighting}[]
\NormalTok{book\_word\_counts }\OtherTok{\textless{}{-}}\NormalTok{ book\_words }\SpecialCharTok{\%\textgreater{}\%}
  \FunctionTok{anti\_join}\NormalTok{(stop\_words) }\SpecialCharTok{\%\textgreater{}\%}
  \FunctionTok{count}\NormalTok{(chapter, word, }\AttributeTok{sort =} \ConstantTok{TRUE}\NormalTok{)}
\end{Highlighting}
\end{Shaded}

\begin{verbatim}
## Joining with `by = join_by(word)`
\end{verbatim}

\begin{Shaded}
\begin{Highlighting}[]
\FunctionTok{head}\NormalTok{(book\_word\_counts)}
\end{Highlighting}
\end{Shaded}

\begin{verbatim}
## # A tibble: 6 x 3
##   chapter word          n
##     <int> <chr>     <int>
## 1      43 elizabeth    36
## 2      18 darcy        33
## 3      18 elizabeth    25
## 4       8 bingley      21
## 5      16 darcy        21
## 6      29 lady         21
\end{verbatim}

Common stopwords such as ``the'', ``and'', and ``of'' are removed
because they do not carry meaningful semantic information. Removing
these words helps the topic model focus on more informative terms. After
removing stopwords, we count how frequently each word appears in each
chapter of the novel.

\begin{Shaded}
\begin{Highlighting}[]
\NormalTok{book\_dtm }\OtherTok{\textless{}{-}}\NormalTok{ book\_word\_counts }\SpecialCharTok{\%\textgreater{}\%}
  \FunctionTok{cast\_dtm}\NormalTok{(chapter, word, n)}

\NormalTok{tm}\SpecialCharTok{::}\FunctionTok{inspect}\NormalTok{(book\_dtm)}
\end{Highlighting}
\end{Shaded}

\begin{verbatim}
## <<DocumentTermMatrix (documents: 61, terms: 6014)>>
## Non-/sparse entries: 26589/340265
## Sparsity           : 93%
## Maximal term length: 17
## Weighting          : term frequency (tf)
## Sample             :
##     Terms
## Docs bennet bingley darcy elizabeth jane lady miss sister time wickham
##   16      3       4    21        18    0    9    5      2    5      17
##   18     10      19    33        25   10   13   11      7   11      16
##   29      2       0     1        12    0   21    9      0    1       1
##   35      1       4     4         2    0    1    2      7    4      11
##   43      0       1    14        36    0    4    7      5   11       4
##   46      0       0     9        11    6    0    1      3    5       5
##   47      7       0     3        18   15    1    3      5    3      11
##   52      2       0    10         5    2    2    2      6    6      14
##   53     12      13     7        14   13    0    3     10    7       2
##   56     18       1     6        19    1   16   15      2    0       0
\end{verbatim}

The dataset is converted into a DocumentTermMatrix (DTM) where each row
represents a chapter and each column represents a word. The values in
the matrix correspond to the number of times a word appears in a
chapter. The resulting matrix contains 61 documents (chapters) and 6014
unique terms. As expected in text data, the matrix is sparse, with about
93\% of the entries being zeros, meaning that most words do not appear
in most chapters.

\begin{Shaded}
\begin{Highlighting}[]
\NormalTok{book\_lda }\OtherTok{\textless{}{-}} \FunctionTok{LDA}\NormalTok{(book\_dtm, }\AttributeTok{k =} \DecValTok{5}\NormalTok{, }\AttributeTok{control =} \FunctionTok{list}\NormalTok{(}\AttributeTok{seed =} \DecValTok{1234}\NormalTok{))}
\NormalTok{book\_topics }\OtherTok{\textless{}{-}} \FunctionTok{tidy}\NormalTok{(book\_lda, }\AttributeTok{matrix =} \StringTok{"beta"}\NormalTok{)}
\FunctionTok{head}\NormalTok{(book\_topics)}
\end{Highlighting}
\end{Shaded}

\begin{verbatim}
## # A tibble: 6 x 3
##   topic term         beta
##   <int> <chr>       <dbl>
## 1     1 elizabeth 0.0172 
## 2     2 elizabeth 0.0125 
## 3     3 elizabeth 0.0127 
## 4     4 elizabeth 0.0223 
## 5     5 elizabeth 0.0146 
## 6     1 darcy     0.00882
\end{verbatim}

We estimate a Latent Dirichlet Allocation (LDA) topic model with five
topics using the DocumentTermMatrix created earlier. LDA identifies
groups of words that frequently appear together across chapters. We then
extract the beta values, which represent the probability that a word
belongs to a specific topic. Higher beta values indicate stronger
associations between words and topics.

\begin{Shaded}
\begin{Highlighting}[]
\NormalTok{book\_top\_terms }\OtherTok{\textless{}{-}}\NormalTok{ book\_topics }\SpecialCharTok{\%\textgreater{}\%}
  \FunctionTok{group\_by}\NormalTok{(topic) }\SpecialCharTok{\%\textgreater{}\%}
  \FunctionTok{slice\_max}\NormalTok{(beta, }\AttributeTok{n =} \DecValTok{10}\NormalTok{) }\SpecialCharTok{\%\textgreater{}\%}
  \FunctionTok{ungroup}\NormalTok{() }\SpecialCharTok{\%\textgreater{}\%}
  \FunctionTok{arrange}\NormalTok{(topic, }\SpecialCharTok{{-}}\NormalTok{beta)}

\NormalTok{book\_top\_terms}
\end{Highlighting}
\end{Shaded}

\begin{verbatim}
## # A tibble: 50 x 3
##    topic term            beta
##    <int> <chr>          <dbl>
##  1     1 elizabeth    0.0172 
##  2     1 jane         0.0131 
##  3     1 bingley      0.0108 
##  4     1 bennet       0.00978
##  5     1 darcy        0.00882
##  6     1 time         0.00769
##  7     1 miss         0.00750
##  8     1 dear         0.00666
##  9     1 sister       0.00635
## 10     1 illustration 0.00499
## # i 40 more rows
\end{verbatim}

To interpret the topics produced by the LDA model, we extract the ten
words with the highest beta values for each topic. These words represent
the most probable terms associated with each topic. The resulting topics
appear to capture major themes and character interactions in the novel.
For example, some topics are dominated by character names such as
Elizabeth, Darcy, and Bennet, reflecting the central relationships in
the story. Other topics include words such as Collins, Lady, and
Catherine, which relate to social hierarchy and marriage proposals.
Another topic includes words such as Wickham and letter, suggesting
narrative moments involving conflict and revelations in the plot.
Overall, the topics capture recurring character interactions and social
dynamics within the novel.

\begin{Shaded}
\begin{Highlighting}[]
\NormalTok{book\_top\_terms }\SpecialCharTok{\%\textgreater{}\%}
  \FunctionTok{mutate}\NormalTok{(}\AttributeTok{term =} \FunctionTok{reorder\_within}\NormalTok{(term, beta, topic)) }\SpecialCharTok{\%\textgreater{}\%}
  \FunctionTok{ggplot}\NormalTok{(}\FunctionTok{aes}\NormalTok{(beta, term, }\AttributeTok{fill =} \FunctionTok{factor}\NormalTok{(topic))) }\SpecialCharTok{+}
  \FunctionTok{geom\_col}\NormalTok{(}\AttributeTok{show.legend =} \ConstantTok{FALSE}\NormalTok{) }\SpecialCharTok{+}
  \FunctionTok{facet\_wrap}\NormalTok{(}\SpecialCharTok{\textasciitilde{}}\NormalTok{ topic, }\AttributeTok{scales =} \StringTok{"free"}\NormalTok{) }\SpecialCharTok{+}
  \FunctionTok{scale\_y\_reordered}\NormalTok{() }\SpecialCharTok{+}
  \FunctionTok{theme\_minimal}\NormalTok{()}
\end{Highlighting}
\end{Shaded}

\pandocbounded{\includegraphics[keepaspectratio]{unsupervised-topicmodels_files/figure-latex/unnamed-chunk-25-1.pdf}}
The figure shows the most important words associated with each of the
five topics estimated by the LDA model. The bar charts display the top
ten words with the highest beta values for each topic, indicating the
strength of association between words and topics. The topics appear to
capture central elements of the novel. Several topics are dominated by
character names such as Elizabeth, Darcy, Jane, and Bennet, reflecting
the importance of interpersonal relationships in the story. One topic
includes words such as Collins, Lady, and Catherine, which relate to
social hierarchy and marriage proposals. Another topic contains words
such as Wickham, Lydia, and letter, suggesting plot developments
involving conflict and revelations. Overall, the topic model
successfully captures recurring character interactions and social
dynamics in the novel.

\begin{Shaded}
\begin{Highlighting}[]
\NormalTok{book\_lda\_8 }\OtherTok{\textless{}{-}} \FunctionTok{LDA}\NormalTok{(book\_dtm, }\AttributeTok{k =} \DecValTok{8}\NormalTok{, }\AttributeTok{control =} \FunctionTok{list}\NormalTok{(}\AttributeTok{seed =} \DecValTok{1234}\NormalTok{))}

\NormalTok{book\_topics\_8 }\OtherTok{\textless{}{-}} \FunctionTok{tidy}\NormalTok{(book\_lda\_8, }\AttributeTok{matrix =} \StringTok{"beta"}\NormalTok{)}

\NormalTok{book\_top\_terms\_8 }\OtherTok{\textless{}{-}}\NormalTok{ book\_topics\_8 }\SpecialCharTok{\%\textgreater{}\%}
  \FunctionTok{group\_by}\NormalTok{(topic) }\SpecialCharTok{\%\textgreater{}\%}
  \FunctionTok{slice\_max}\NormalTok{(beta, }\AttributeTok{n =} \DecValTok{10}\NormalTok{) }\SpecialCharTok{\%\textgreater{}\%}
  \FunctionTok{ungroup}\NormalTok{() }\SpecialCharTok{\%\textgreater{}\%}
  \FunctionTok{arrange}\NormalTok{(topic, }\SpecialCharTok{{-}}\NormalTok{beta)}
\end{Highlighting}
\end{Shaded}

\begin{Shaded}
\begin{Highlighting}[]
\NormalTok{book\_top\_terms\_8 }\SpecialCharTok{\%\textgreater{}\%}
  \FunctionTok{mutate}\NormalTok{(}\AttributeTok{term =} \FunctionTok{reorder\_within}\NormalTok{(term, beta, topic)) }\SpecialCharTok{\%\textgreater{}\%}
  \FunctionTok{ggplot}\NormalTok{(}\FunctionTok{aes}\NormalTok{(beta, term, }\AttributeTok{fill =} \FunctionTok{factor}\NormalTok{(topic))) }\SpecialCharTok{+}
  \FunctionTok{geom\_col}\NormalTok{(}\AttributeTok{show.legend =} \ConstantTok{FALSE}\NormalTok{) }\SpecialCharTok{+}
  \FunctionTok{facet\_wrap}\NormalTok{(}\SpecialCharTok{\textasciitilde{}}\NormalTok{ topic, }\AttributeTok{scales =} \StringTok{"free"}\NormalTok{) }\SpecialCharTok{+}
  \FunctionTok{scale\_y\_reordered}\NormalTok{() }\SpecialCharTok{+}
  \FunctionTok{theme\_minimal}\NormalTok{()}
\end{Highlighting}
\end{Shaded}

\pandocbounded{\includegraphics[keepaspectratio]{unsupervised-topicmodels_files/figure-latex/unnamed-chunk-27-1.pdf}}
This figure presents the topics identified when the number of topics is
increased to eight. Compared to the five-topic model, the eight-topic
model produces more fine-grained distinctions between themes in the
novel. While several topics continue to be dominated by character names
such as Elizabeth, Darcy, and Bennet, additional topics capture more
specific narrative elements, such as interactions involving Mr Collins,
the Wickham storyline, and family relationships. Increasing the number
of topics allows the model to separate certain narrative elements into
more specialized clusters. However, because character names dominate the
vocabulary of the novel, some overlap between topics remains. Overall,
the eight-topic model provides more detailed thematic distinctions but
may be slightly harder to interpret compared to the five-topic model.

Overall, the topic models capture key themes in the novel, particularly
interpersonal relationships and social dynamics among central
characters. The five-topic model provides a clearer and more
interpretable overview of the main narrative themes, while the
eight-topic model reveals more specific sub-themes within the story.
This comparison highlights how the choice of the number of topics
influences the granularity and interpretability of topic modeling
results.

\end{document}
